%%%% CAPÍTULO 5 - CONCLUSÕES E PERSPECTIVAS

\chapter{Conclusão}\label{cap:conclusoeseperspectivas}

Este trabalho apresentou o desenvolvimento e a documentação de uma plataforma de HFT de baixo custo baseada em co-projeto ARM--FPGA, com foco principal na criação de estratégias de trading em MATLAB e geração de hardware por meio do MATLAB HDL Coder. A infraestrutura desenvolvida cobre o caminho necessário para alimentar uma estratégia gerada por HLS: dados sintéticos de mercado, normalização no ARM, comunicação com a FPGA, livro de ofertas simplificado, decisão em hardware, validação por testes e preparação para execução na DE10-Nano.

O objetivo geral foi atendido no nível funcional: o sistema permite isolar a estratégia como função MATLAB, gerar VHDL com HDL Coder, integrar o bloco gerado ao motor de decisão e executar o caminho de dados entre software e FPGA. O trabalho também consolidou o contrato entre C++ e VHDL, enfatizando registradores, ponteiros, capacidade efetiva de filas, cálculo dinâmico da base RX, largura de oito palavras por quadro e ordem de empacotamento. Esse contrato é uma contribuição de engenharia essencial, pois mantém estável a infraestrutura sobre a qual novas estratégias podem ser testadas.

Entre as vantagens observadas, destacam-se a modularidade da arquitetura, a separação adequada entre software flexível e hardware determinístico, a possibilidade de alterar a regra de estratégia em MATLAB e regenerá-la por HDL Coder, e a existência de testes em diferentes níveis. A abordagem por filas circulares é simples, dispensa bloqueios entre produtor e consumidor e facilita a depuração. A modelagem do livro por profundidade fixa favorece síntese em FPGA e evita estruturas dinâmicas incompatíveis com RTL.

Os objetivos específicos também foram contemplados. A plataforma foi organizada para que a estratégia seja o principal bloco substituível pelo usuário. O estudo teórico apresentou os conceitos necessários de HFT, livro de ofertas, FPGA, HLS, MATLAB HDL Coder, FAST, MMIO, Avalon-MM e métricas de desempenho. A documentação do contrato ARM--FPGA foi incorporada ao texto e detalhada em apêndice. Os módulos de software, hardware e infraestrutura foram descritos com suas funções, vantagens, limitações e integração. Por fim, a validação funcional foi organizada em uma matriz que relaciona testes a riscos técnicos.

Do ponto de vista acadêmico, a contribuição central não é uma estratégia de negociação lucrativa, mas um fluxo de criação de estratégias para FPGA. Essa distinção é importante. Estratégias de HFT reais dependem de dados proprietários, controle de risco, colocação de ordens, posição, custos e conformidade. A plataforma, por sua vez, demonstra como uma aplicação sensível a latência pode ser decomposta em software e hardware, como representar eventos em ponto fixo, como preservar ordenação em uma fila circular e como usar HLS para reduzir o esforço de implementação do bloco decisório.

As limitações também são claras. O feed é sintético e usa TCP, não dados reais de bolsa. O livro de ofertas é por nível agregado, não por ordem individual. A estratégia é didática e não considera posição, risco, custos, execução ou conformidade regulatória. As medições quantitativas de latência, vazão e uso de recursos ainda precisam ser realizadas com instrumentação adequada. Além disso, MMIO com polling é suficiente para o protótipo desenvolvido, mas pode não ser adequado para taxas profissionais de dados de mercado.

Outra limitação é o estágio das métricas. A execução funcional em placa confirma que o caminho feed--ARM--FPGA--ARM opera, mas a avaliação de desempenho ainda precisa de instrumentação consistente. Para reivindicar desempenho em placa, é necessário coletar timestamps, contadores de ciclo, vazão máxima, uso de recursos e frequência de operação. Essa cautela evita transformar uma validação funcional em afirmação quantitativa não comprovada.

Os resultados quantitativos obtidos com o benchmark revelam uma distinção importante entre aceleração de núcleo lógico e latência fim a fim. A FPGA executa o núcleo de livro e decisão em 60~ns com jitter zero, mas o custo de comunicação pela ponte MMIO eleva o RTT ARM$\to$FPGA$\to$ARM para cerca de 9.700~ns. Esse resultado, combinado com o software C++ equivalente em 311~ns, mostra que o gargalo atual não está na lógica de hardware, mas na interface de comunicação. Sistemas de produção resolvem esse problema posicionando a FPGA diretamente no caminho de rede, mas isso ultrapassa o escopo deste protótipo.

Como trabalhos futuros, recomenda-se:

\begin{itemize}
  \item \textbf{Transporte por DMA}: substituir o acesso MMIO word a word por transferências DMA (por exemplo, FPGA Manager DMA ou AXI DMA no Cyclone~V), reduzindo o overhead de comunicação de $\sim$9.700~ns para $\sim$500--1.000~ns por mensagem. Isso aproximaria a latência fim a fim da latência lógica interna.

  \item \textbf{FPGA no caminho de dados de rede}: integrar recepção Ethernet diretamente na FPGA (via SGMII ou interface TSE), eliminando o ARM do laço crítico. Nessa arquitetura, o fluxo seria Ethernet $\to$ parser FPGA $\to$ livro + estratégia $\to$ envio de ordem $\to$ Ethernet, com latência de fim a fim na faixa de 100--500~ns conforme reportado na literatura \cite{Leber2011,Lockwood2012}.

  \item \textbf{Medição de energia por decisão}: instrumentar o consumo de corrente do Cyclone~V com medição direta (amperímetro no trilho de alimentação) durante o benchmark, obtendo energia por decisão empírica em vez da estimativa por ordem de grandeza apresentada na \autoref{tab:energiaComparacao}.

  \item \textbf{Coleta de relatórios Quartus}: coletar os relatórios de síntese e temporização do Quartus Prime (elementos lógicos, registradores, blocos M10K, multiplicadores e frequência máxima por módulo) e substituir o \autoref{gra:placeholderRecursos} por dados reais.

  \item \textbf{Estratégias mais complexas em MATLAB}: avaliar regras com médias móveis, bandas de Bollinger ou correlação entre símbolos, medindo o impacto no uso de recursos e na latência lógica à medida que a complexidade do código MATLAB aumenta.

  \item \textbf{Modelo de eventos enriquecido}: ampliar para cancelamentos, execuções parciais e múltiplos níveis de profundidade, validando a robustez da infraestrutura frente a padrões de livro mais próximos do real.

  \item \textbf{Controle de risco e posição}: adicionar limites de posição, janela de perda máxima e conformidade de ordens como módulo separado, mantendo a estratégia MATLAB como bloco isolado da lógica de risco.
\end{itemize}

A evolução natural do projeto pode ser dividida em três frentes. A primeira é experimental: coletar métricas com instrumentação adequada (relatórios Quartus, timestamps por etapa, medição de energia). A segunda é arquitetural: substituir MMIO por DMA ou por recepção de rede direta em FPGA, eliminando o gargalo de comunicação identificado neste trabalho. A terceira é financeira: enriquecer o modelo de eventos, incluir controle de posição, limites de risco e custos de transação. Essas frentes são complementares e devem ser conduzidas nessa ordem para evitar que uma estratégia mais complexa esconda problemas básicos de infraestrutura.

Conclui-se que a proposta de uma plataforma de HFT de baixo custo orientada à geração de estratégias por MATLAB HDL Coder é viável no nível funcional apresentado. Mesmo sem atingir o escopo de um sistema financeiro real, o protótipo demonstra que uma estratégia escrita em MATLAB pode ser tratada como bloco substituível, gerada para VHDL por HLS e integrada a uma infraestrutura ARM--FPGA com livro de ofertas, filas MMIO e validação cruzada. Esses elementos constituem uma base sólida para a continuidade do projeto, especialmente na coleta de métricas quantitativas e na avaliação de novas estratégias descritas em MATLAB.
